% Generated draft: requires mathematical and source review.
\documentclass[11pt]{article}
\usepackage[margin=25mm]{geometry}
\usepackage{fontspec}
\setmainfont{Latin Modern Roman}
\usepackage{amsmath,amssymb,mathtools,unicode-math}
\setmathfont{Latin Modern Math}
\usepackage{enumitem,needspace,xurl,hyperref}
\setlength{\parindent}{0pt}
\setlength{\parskip}{5pt}
\setlength{\emergencystretch}{3em}
\newcommand{\N}{\mathbb{N}}
\newcommand{\Z}{\mathbb{Z}}
\newcommand{\Q}{\mathbb{Q}}
\newcommand{\R}{\mathbb{R}}
\newcommand{\C}{\mathbb{C}}
\newcommand{\F}{\mathbb{F}}
\newcommand{\A}{\mathbb{A}}
\newcommand{\PP}{\mathbb{P}}
\newcommand{\E}{\mathbb{E}}
\newcommand{\eps}{\varepsilon}
\newcommand{\dd}{\,\mathrm{d}}
\begin{document}
\section*{UnsolvedMath ID: 30004394}
\textbf{Source problem code:} catalogue entry
\subsection{Definitions and mathematical statement}
Let \(d\in\Z\) with \(d\geq3\), and let \((M_n)_{n\geq1}\) be triangulated compact \(d\)-manifolds, with triangulations \(K_n\). Write \(V(K_n)\) for the vertex set and \(K_n^{(1)}\) for the graph consisting of vertices and edges. The uniform degree assumption is that there exists \(D<\infty\) such that
\[
\deg_{K_n^{(1)}}(v)\leq D
\quad\text{for every }n\geq1\text{ and every }v\in V(K_n).
\]
The source asks whether Benjamini--Schramm convergence to a triangulation of \(\R^d\) implies
\[
\lim_{n\to\infty}\frac{\dim H_1(M_n)}{|M_n|}=0.
\]
Here the coefficient system for homology, the meaning of \(|M_n|\), the treatment of manifold boundary, and the precise rooted interpretation of the limit are not specified in the supplied record. Consequently, this displayed question does not yet determine a unique mathematical assertion.

One explicit conditional formulation, if finite triangulations, vertex normalization, homology over a fixed field, and simplicial Benjamini--Schramm convergence are intended, is as follows. Fix a field \(\Bbbk\) and put
\[
 b_1(M_n;\Bbbk)=\dim_{\Bbbk}H_1(M_n;\Bbbk).
\]
Choose a root \(o_n\) uniformly from the nonempty finite set \(V(K_n)\). For a rooted simplicial complex \((K,o)\), let \(B_r(K,o)\) be the rooted subcomplex induced by the vertices at graph distance at most \(r\) from \(o\) in \(K^{(1)}\). Simplicial Benjamini--Schramm convergence to a random rooted locally finite complex \((T,o)\) means that, for every integer \(r\geq0\) and every finite rooted simplicial complex \((A,a)\),
\[
\lim_{n\to\infty}\Pr\bigl(B_r(K_n,o_n)\cong(A,a)\bigr)
=\Pr\bigl(B_r(T,o)\cong(A,a)\bigr),
\]
where isomorphisms preserve the root and all simplices. Under the interpretation that \(T\) triangulates \(\R^d\) almost surely, the conditional question is whether every such sequence satisfies
\[
\lim_{n\to\infty}\frac{b_1(M_n;\Bbbk)}{|V(K_n)|}=0.
\]
This conditional formulation supplies possible meanings for missing notation; it is not asserted to be the uniquely intended source formulation. In particular, neither the choice of \(\Bbbk\) nor universal quantification over fields is supplied. The source also does not specify whether its limiting triangulation is fixed or random, whether convergence retains all simplices or only the \(1\)-skeleton, or whether compact manifolds with boundary are admitted. These qualifications remain unresolved; no connectedness or orientability hypothesis is supplied.
\subsection{Short English statement}
The question asks whether compact triangulated manifolds of dimension at least three, with uniformly bounded vertex degrees and local limit a triangulation of \(\R^d\), must have vanishing normalized first homology. A precise assertion requires clarification of the homology coefficients, normalization, boundary convention, and local convergence model. The supplied status is \emph{partially solved}; this draft does not independently verify that assessment.
\subsection{Sources}
\begin{itemize}
\item Local UnsolvedMath registry: supplied complete record for permanent ID \(30004394\), titled \emph{Vanishing Normalized Homology Under Local Manifold Convergence}, in the Oberwolfach Reports Open Problems collection. Exact source problem code: \texttt{catalogue entry}. The proposed year is recorded as 2020; no proposer is identified.
\item Supplied source citation: \emph{Manifolds and Groups (2021)}, Oberwolfach Reports, DOI \texttt{10.4171/owr/2020/7}, pp.~417--463. \url{https://doi.org/10.4171/owr/2020/7}. The supplied background describes Oberwolfach Reports as documenting mathematical workshops and their open-problem discussions.
\item Supplied literature reference: Hutchcroft--Pete, \emph{Kazhdan groups have cost 1}. \url{https://arxiv.org/abs/1810.11015}. The record does not identify a particular theorem establishing the stated manifold consequence.
\item Supplied dated literature curation, checked at \texttt{2026-08-22}: status \emph{partially solved}. The supplied assessment says that the Benjamini--Schramm first-homology limit is known in dimension three, for covering towers in special settings, and under a Kazhdan-property hypothesis, while the general question remains open and includes the degree-one Singer problem. This is supplied dated curation derived from the OWR material, not an independently verified literature assessment. The supplied background expressly requests independent verification of time-sensitive status and cited claims.
\end{itemize}
\end{document}
